\section*{第 \ref{ch:g4:breathing} 章——呼吸}

\begin{solution}{exo:g4:breathing:1}
\omterm{def:g4:breathing:movements}{吸气}：胸腔扩张，空气流入。\omterm{def:g4:breathing:movements}{呼气}：
胸腔放松，空气流出。
\end{solution}

\begin{solution}{exo:g4:breathing:2}
鼻（或口）、\omterm{def:g4:breathing:organs}{气管}、两个\omterm{def:g1:plants-around-us:tree}{分支}、肺，
最后是数百万个微小气囊，
外面包着血管。
\end{solution}

\begin{solution}{exo:g4:breathing:3}
约每分钟 $15$ 次。冲刺时可升至 $40$
次或更多，\omterm{def:g4:breathing:movements}{呼吸}也更深——工作中的\omterm{def:g3:bones-and-muscles:muscle}{肌肉}
需要更快获得氧气，因此肺加快吸入氧气。
\end{solution}

\begin{solution}{exo:g4:breathing:4}
它吸入氧气，并排出二氧化碳。
这种交换发生在肺的微小气囊中，
穿过薄壁，与\omterm{prop:g4:journey-of-food:absorb}{血液}进行。
\end{solution}

\begin{solution}{exo:g4:breathing:5}
呼出的空气更温暖、更湿润（能使镜子蒙上水雾），
氧气更少，二氧化碳更多。
\end{solution}

\begin{solution}{exo:g4:breathing:6}
因为空气与\omterm{prop:g4:journey-of-food:absorb}{血液}的交换需要表面积：数百万个气囊壁
的面积相加，可达半个网球场，折叠在胸腔内。
\omterm{def:g4:journey-of-food:tract}{小肠}也用类似的办法，借助其皱褶
吸收\omterm{def:g4:journey-of-food:digestion}{营养物质}。
\end{solution}

\begin{solution}{exo:g4:breathing:7}
鼻子在空气进入时对其加温、加湿并过滤——
寒冷、干燥、多尘的空气，
到达娇嫩的肺时已经处理妥当。
\end{solution}

\begin{solution}{exo:g4:breathing:8}
开放性答案。静息时通常约为 $15$ 次，
做完开合跳后为 $25$--$40$ 次：工作中的\omterm{def:g3:bones-and-muscles:muscle}{肌肉}需要更多氧气，
因此\omterm{def:g4:breathing:movements}{呼吸}加快、加深，以供应所需。
\end{solution}

\begin{solution}{exo:g4:breathing:9}
人体储存的食物很多（可供数小时或数天之用），
但几乎不储存氧气——仅够一分钟左右。
因此氧气必须不断输送，一次\omterm{def:g4:breathing:movements}{呼吸}接着一次\omterm{def:g4:breathing:movements}{呼吸}，
这就是\omterm{def:g4:breathing:movements}{呼吸}绝不能长时间停顿的原因。
\end{solution}

\begin{solution}{exo:g4:breathing:10}
二十五人整个上午都在吸入房间里的氧气、呼出二氧化碳，
房间空气因此氧气减少、二氧化碳增多。
解决办法是第 1 条规则：
打开窗户，通风换气。
\end{solution}

\begin{solution}{exo:g4:breathing:11}
两者都是\omterm{prop:g4:journey-of-food:absorb}{血液}与氧气相遇的场所：
巨大、薄而布满血管的表面，氧气由此进入，
二氧化碳由此排出——鳃从水中获取溶解的氧气，
肺则从空气中获取氧气。在空气中，虹鳟的
\omterm{def:g1:animals-around-us:covering}{羽毛}状鳃表面会塌陷并粘连在一起：巨大的交换面积
几乎缩小为零，没有交换面积，
无论周围空气多么充足，也得不到氧气。
\end{solution}
